% !TeX program = xelatex
% UTF-8; faithful transcription of nine handwritten Persian/English sheets.
% Source: Google Drive / Notebook / congenital infections.
% Clinical statements and numerical values are preserved as written, not corrected.
% pic01.jpg = 20261016_114502288.jpg (1BQNQV3HboVh0XMlziRBzFvq45ONaPHM1)
% pic02.jpg = 20261016_114507766.jpg (1PEWVlesxtkT8xkEN-rruxydokivVXM-u)
% pic03.jpg = 20261016_114512822.jpg (13cCw7fu2nO409VW0PFmpTPhykYO57H_V)
% pic04.jpg = 20261016_114518415.jpg (12jmGqdbztHuZ1kX6w4LlzUzw42d_2zip)
% pic05.jpg = 20261016_114522866.jpg (1LYdmhymXr12lzCtaQ8tpOoxf9aMTDBWG)
% pic06.jpg = 20261016_114528206.jpg (1gziHtCbsfgy9UxWsBpKFb5ekpSJm3lAD)
% pic07.jpg = 20261016_114532757.jpg (1A1TLSwihE4qpvIbYDnYOTWmfrMROva5w)
% pic08.jpg = 20261016_114537767.jpg (1Xj7v9zzpaNLH_Z2jexrpJSxF1li3Y5Pq)
% pic09.jpg = 20261016_114540949.jpg (1kcATN4m2fxNAMookCmqarBOmzOWteCOk)
\documentclass[11pt,a4paper]{article}
\usepackage[margin=18mm]{geometry}
\usepackage{amsmath,amssymb}
\usepackage{enumitem}
\usepackage{xcolor}
\usepackage[normalem]{ulem}
\usepackage{xepersian}
\settextfont{Amiri}
\setlatintextfont{DejaVu Serif}
\setlength{\parindent}{0pt}
\setlength{\parskip}{4pt}
\setlength{\emergencystretch}{3em}
\setlist[itemize]{itemsep=4pt,topsep=4pt,leftmargin=*}
\newcommand{\en}[1]{\lr{#1}}
\newcommand{\hl}[1]{\begingroup\setlength{\fboxsep}{1pt}\colorbox{yellow!40}{\strut #1}\endgroup}
\newcommand{\unclear}[1]{\textcolor{red!65!black}{[خوانش نامطمئن: #1]}}
\newcommand{\sheet}[1]{\section*{تصویر #1}\begingroup\small
سربرگ: \en{Subject:} ۹۶ ـ عفونت‌های همراه با تولد؛
\en{Year:}\quad\en{Month:}\quad\en{Date:}؛ شمارهٔ پایین برگه: #1.\par\endgroup}
\begin{document}
\begin{center}{\Large\bfseries عفونت‌های همراه با تولد}\end{center}
\sheet{۱}
\begin{itemize}
\item عفونت همراه با تولد: \hl{طی بارداری از طریق جفت} به جنین انتقال می‌یابد.
\item پاتوژن‌های شایع عفونت‌های همراه با تولد: توکسوپلاسما، روبلا، \en{CMV}، \en{HSV}، \en{VZV}، سیفلیس مادرزادی، پارووویروس، \en{HIV}، \en{HBV}، نایسریا گنوره، زیکا ویروس، کلامیدیا، \en{TB}.
\item تظاهرات بالینی این عفونت‌ها بسیار شبیه به هم است و شامل:
\en{IUGR}، هیدروپس غیرایمن، آنمی، ترومبوسایتوپنی، زردی، ارگانومگالی، کوریورتینیت، مالفورماسیون‌ها.
\item ارزیابی نوزاد با شک به عفونت همراه با تولد:
\begin{itemize}
\item ایزوله کردن پاتوژن؛ کشت: روبلا، \en{CMV}، \en{HSV}، گنوره، \en{TB}.
\item یافتن \en{Ag}: کلامیدیا، \en{HBV}.
\item \en{PCR}.
\item یافتن \en{Ab}: \en{IgG} برای توکسوپلاسما، سیفلیس، پارووویروس، \en{HIV}، بورلیا.
\end{itemize}
\end{itemize}
\subsection*{توکسوپلاسما}
\begin{itemize}
\item عفونت حاد مادر، سپس انتقال به جنین از راه جفت؛ \hl{ریسک ۴۰\%}.
\item اگر بیماری در بدن مادر وجود داشته باشد، سپس فعال شود، به جنین \hl{انتقال نمی‌یابد}.
\item ریسک انتقال ویروس به جنین \hl{وابسته به سن حاملگی و میزان جریان خون جفت} است.
\begin{itemize}
\item در \hl{جنین ترم} ریسک انتقال نزدیک به ۹۰\% است.
\end{itemize}
\item اکثر نوزادان علائم تحت بالینی دارند و ما حتی متوجه بیماری نمی‌شویم.
\item علائم کلاسیک: \hl{هیدروسفالی}، \hl{کوریورتینیت}، \hl{اینتراکرانیال کلسیفیکیشن}.
\end{itemize}
\clearpage
\sheet{۲}
\begin{itemize}
\item نوزادان بیمار نسبتاً جثه‌شان کوچک هستند.
\begin{itemize}
\item دچار زردی زودرس می‌شوند.
\item دارای \hl{ارگانومگالی} هستند.
\item \hl{راش‌های ماکولوپاپولر جنرالیزه}.
\item \hl{تشنج} تظاهر شایعی است.
\end{itemize}
\item در تصویربرداری از مغز: کلسیفیکاسیون منتشر کورتیکال؛ در \en{CMV}، \hl{اکثراً پری‌ونتریکولار} است.
\item نوزادان مبتلا، ریسک افزایش‌یافتهٔ عوارض نورودولوپمنتال تأخیری دارند.
\item تشخیص توکسوپلاسموز در نوزاد:
\begin{itemize}
\item \hl{\en{Serum IgG Ab}}.
\item ۱ تا ۲ ماه پس از ابتلا به حداکثر میزان می‌رسد.
\item اگر \en{IgG Ab} حداقل ۴ برابر \en{N/L} بود، برای تأیید تشخیص \hl{\en{IgM Ab}} را درخواست می‌کنیم.
\item \en{IgE} و \en{IgA} هم می‌توانند کمک کنند.
\end{itemize}
\item \hl{بررسی چشمی، شنوایی و نورولوژیک} (\en{CT}، \en{LP}): \hl{الزامی است}.
\item درمان: هم باعلامت، هم بدون علامت را درمان می‌کنیم.
\begin{itemize}
\item \hl{پریمتامین + اسید فولینیک + سولفادیازین}.
\end{itemize}
\item طول درمان: تا یک سال.
\end{itemize}
\subsection*{روبلا}
\begin{itemize}
\item در صورت ابتلا در \hl{چهار هفتهٔ} اول، احتمال آنومالی ۸۵\% است.
{\small (عبارت زیر اصلاح «چهار هفته»: \sout{سه‌ماهه}.)}
\item به ۴۰\% موارد می‌رسد.
\item \hl{ابتلا پس از ماه ۴ حاملگی معمولاً نقصی ایجاد نمی‌کند}.
\end{itemize}
\clearpage
\sheet{۳}
\subsection*{روبلا ـ ادامه}
\begin{itemize}
\item شایع‌ترین تظاهرات:
\begin{itemize}
\item چشمی: کاتاراکت، رتینوپاتی، گلوکوم.
\item قلبی: \hl{\en{PDA}}، \hl{استنوز پریفرال شریان پولی}.
\item شنوایی: کری حسی عصبی.
\item نورولوژیک: مننگوانسفالیت، اختلالات رفتاری، تأخیر رشد.
\end{itemize}
\item دیگر تظاهرات:
\begin{itemize}
\item عقب‌ماندگی رشد جسمی.
\item \hl{لنفادنوپاتی}.
\item زردی زودرس.
\item \hl{ترومبوسایتوپنی}.
\item درمال اریتروپوئز: \hl{بلوبری مافین}.
\item \en{Radiolucent bone disease}.
\end{itemize}
\item تشخیص: \hl{\en{Rubella IgM Ab}}؛ تأیید تشخیص: اندازه‌گیری \hl{\en{Rubella IgG Ab} طی چند ماه}.
\item ویروس روبلا در خون، ادرار، \en{CSF} و مخاط گلو قابل شناسایی است.
\item نوزاد مبتلا به روبلا تا یک سال ویروس را به مایعات بدن دفع می‌کند؛ مدفوع.
\item نوزاد در بیمارستان ایزوله شود؛ به‌خصوص از زنان باردار و افراد مستعد جدا شود.
\end{itemize}
\subsection*{\en{CMV}}
\begin{itemize}
\item \hl{شایع‌ترین علت عفونت همراه با تولد} است.
\item عامل اول \hl{کری حسی عصبی}، \hl{\en{CP}}، عقب‌ماندگی ذهنی و بیماری شبکیه.
\item اگر مادر باردار مبتلا به \en{CMV} شود، احتمال انتقال به جنین ۳۵\% است؛ مادران نوجوان با احتمال بیشتری منتقل می‌کنند.
\item \hl{ریسک انتقال به جنین مستقل از سن حاملگی است}.
\end{itemize}
\clearpage
\sheet{۴}
\subsection*{\en{CMV} ـ ادامه}
\begin{itemize}
\item هرچه مادر در سن بارداری پایین‌تری مبتلا شود، علامت‌های نوزاد شدیدتر خواهد بود.
\item شایع‌ترین مسیرهای انتقال ویروس به مادر باردار: ارتباط جنسی، اشتغال نزد کودکان.
\item حتی اگر مادر قبل از حاملگی نیز مبتلا به \en{CMV} \unclear{بشوده / نبوده} می‌تواند آن را انتقال دهد.
\item بیش از ۹۰\% نوزادانی که در جنینی مبتلا به \en{CMV} شده‌اند هیچ علامتی ندارند.
\item یافته‌ها: \hl{میکروسفالی}، \hl{ترومبوسایتوپنی}، \hl{ارگانومگالی}، هپاتیت، \hl{کلسیفیکاسیون داخل جمجمه}، کوریورتینیت، اختلال شنوایی.
\item در برخی نوزادان \hl{بلوبری مافین} به پی اریتروپوئز در پوست.
\item برخی علایم تا چند سال بعد مشخص نمی‌شوند: کاهش شنوایی، عقب‌ماندگی رشد.
\item ابتلا به \en{CMV} حین زایمان و یا از طریق شیر منجر به بیماری و آسیب \en{CNS} نمی‌شود.
\item تشخیص:
\begin{itemize}
\item یافتن ویروس در ادرار یا بزاق.
\item \hl{یافتن ویروس طی ۳ هفتهٔ اول تولد}: تشخیص عفونت مادرزادی.
\end{itemize}
\item درمان: \hl{گانسیکلوویر}؛ عوارض شنوایی را کم می‌کند.
\end{itemize}
\subsection*{\en{HSV}}
\begin{itemize}
\item عامل ۸۰\% عفونت‌های هرپسی نوزادان \en{HSV-2} است.
\item ابتلا به \en{HSV} طی دورهٔ جنینی، \hl{ریسک زایمان زودرس} را بالا می‌برد.
\item اگر مادر مبتلا باشد و زایمان طبیعی انجام دهد، با احتمال حدود ۵۰\% نوزاد مبتلا می‌شود.
\end{itemize}
\clearpage
\sheet{۵}
\subsection*{\en{HSV} ـ ادامه}
\begin{itemize}
\item علائم عفونت ظرف \hl{۵ تا ۷ روز} پس از تولد ظاهر می‌شود (تا ۴ هفتهٔ اول).
\item علائم: درگیری ارگان‌های مختلف، عمدتاً کبد و ریه؛ \en{CNS}، پوست، چشم و دهان.
\item به چه کسانی باید شک کرد؟
\begin{itemize}
\item هر نوزاد با علائم سپسیس، اختلال عملکرد کبد و کشت منفی برای باکتری.
\item \hl{هر نوزادی با تب، بی‌قراری، \en{LP} غیرطبیعی و تشنج}.
\end{itemize}
\item معمولاً بیماری شدیدی ایجاد می‌کند و اگر درمان نشود، عوارض جدی به جا می‌گذارد.
\item تشخیص: نمونه‌برداری از وزیکول‌های پوستی، نازوفارنکس، چشم، خون، ادرار، مدفوع یا رکتوم، \en{CSF}.
\item \hl{کشت مثبت ظرف ۴۸ ساعت پس از تولد}: تأیید ابتلا حین دورهٔ جنینی.
\item از \en{PCR} نیز می‌توان برای خون، \en{CSF} و وزیکول‌ها استفاده کرد.
\item درمان: \hl{آسیکلوویر وریدی}؛ حتی به صورت شک به \en{HSV} درمان را شروع می‌کنیم.
\end{itemize}
\subsection*{سیفلیس همراه با تولد}
\begin{itemize}
\item ابتلا عمدتاً از طریق جفت است، سپس کانال زایمانی.
\item عوارض ابتلا به سیفلیس در جنینی: سقط، مرده‌زایی، هیدروپس، پره‌مچوریتی.
\item علائم: \hl{ارگانومگالی}، \hl{لنفادنوپاتی}، \hl{ترومبوسیتوپنی}، \hl{آنمی همولیتیک}، راش پوستی، استئوکندریت، ضایعات خاص عصبی، \en{snuffles}.
\end{itemize}
\clearpage
\sheet{۶}
\subsection*{سیفلیس همراه با تولد ـ ادامه}
\begin{itemize}
\item اگر درمان نشوند، ممکن است علائم را پس از \hl{۲ سالگی} نشان دهند:
درگیری \en{CNS}، استخوان، دندان، غضروف، چشم، پوست.
\item علائم تأخیری: \hl{دندان‌های هاچینسون}، کراتیت بینابینی، \hl{کری ناشی از عصب ۸}،
\en{bowing of the shins}، \en{frontal bossing}، \en{mulberry molars}، \en{rhagades}، \hl{\en{Clutton joints}}.
\item \hl{تریاد هاچینسون}: درگیری دندان‌ها، کراتیت بینابینی، کری عصبی.
\item اکثراً در زمان تشخیص بی‌علامت؛ ۲ تا ۵ هفتگی علامت‌دار می‌شوند.
\item علائم اولیه: \hl{\en{snuffles + poor feeding}} (رینیت سیفلیسی).
در مقایسه با سرماخوردگی طولانی‌تر، شدیدتر و معمولاً \hl{خونی} هستند.
\item \hl{راش‌های ماکولوپاپولر همراه با پوسته‌ریزی}؛ \hl{کف دست و پا و اطراف دهان و مقعد}.
\item می‌تواند باعث \hl{نوریت} و \hl{کوریورتینیت} شود.
\item بررسی مایع \en{CSF}: پلئوسیتوز، افزایش پروتئین.
\item یافته‌های بسیار شایع: درگیری استخوان‌های دراز + استئوکندریت و پریوستیت.
\item هر زنی که در بیمارستان زایمان می‌کند باید سرولوژی سیفلیس بررسی شود.
\item اگر مادر مثبت شد، نوزادش باید بررسی شود: \en{VDRL}، \en{RPR}.
\begin{itemize}
\item \en{RPR}: \en{Rapid Plasma Reagin}.
\end{itemize}
\end{itemize}
\clearpage
\sheet{۷}
\subsection*{سیفلیس همراه با تولد ـ ادامه}
\begin{itemize}
\item تست مادر و نوزاد باید از یک نوع باشد.
\item اندیکاسیون‌های بررسی بیشتر و درمان نوزاد:
\begin{itemize}
\item تیتر مادر در دو بررسی \hl{بیش از ۴ برابر} شود.
\item تیتر نوزاد \hl{بیش از ۴ برابر مادر} باشد.
\item \hl{نوزاد علامت‌دار} شود.
\item \hl{مادر درمان درست برای سیفلیس نگرفته باشد}.
\end{itemize}
\item نوزاد ممکن است مبتلا باشد ولی سرولوژی منفی باشد.
\item \en{CSF} هم باید \en{VDRL} و هم از نظر سلول‌شمار و پروتئین بررسی شود.
\begin{itemize}
\item \hl{\en{VDRL} مثبت تشخیصی است}؛ \en{WBC} و \en{pro} افزایش‌یافته است.
\end{itemize}
\item درمان: پنی‌سیلین وریدی؛ \hl{\en{Penicillin G for 10--14 days}}.
\item باید پس از ۳، ۶ و ۱۲ ماه از درمان آزمایش برای تأیید کاهش \en{Ab} انجام شود.
\item اگر نوروسیفلیس وجود داشت، \hl{تا ۳ سال هر ۶ ماه باید \en{LP}} انجام شود؛ یا تا زمانی که \en{CSF} نرمال گردد.
\end{itemize}
\subsection*{نایسریا گنوره}
\begin{itemize}
\item معمولاً چشم را درگیر می‌کند.
\item باکتریمی، آرتریت، مننژیت، واژینیت، آبسهٔ اسکالپ.
\item انتقال به نوزاد از طریق کانال زایمان.
\item علائم ظرف \hl{۵ روز اول} بروز می‌کند؛ ترشح آبکی چشم که \hl{چرکی} می‌شود.
\item درگیری چشم‌ها معمولاً \hl{دوطرفه} است.
\end{itemize}
\clearpage
\sheet{۸}
\subsection*{نایسریا گنوره ـ ادامه}
\begin{itemize}
\item اگر درمان نشود، قرنیه و اتاقک قدامی را درگیر می‌کند؛ کوری.
\item درمان: \hl{سفتریاکسون عضلانی تک‌دوز}؛ شستشوی چشم با محلول سالین.
\item \hl{آنتی‌بیوتیک موضعی جایگاهی ندارد}.
\item نوزادان با درگیری چشمی باید بستری شوند و از نظر درگیری منتشر بررسی شوند: سپس، آرتریت، مننژیت؛ \hl{درمان با سفتریاکسون} به \hl{مدت ۷ روز}.
\item اگر نوزاد هایپربیلی‌روبینمی داشت: درمان با \hl{سفپیم}.
\item نوزاد با مننژیت: درمان ۱۰ تا ۱۴ روز.
\item بیمار از نظر کلامیدیا، سیفلیس و \en{HIV} باید بررسی شود.
\item پروفیلاکسی برای هر نوزادی: سیلور نیترات، \hl{اریترومایسین} یا تتراسایکلین.
\end{itemize}
\subsection*{کلامیدیا}
\begin{itemize}
\item انتقال از طریق کانال زایمان؛ احتمال انتقال ۵۰\% در پی زایمان طبیعی.
\item بیشترین محل درگیری: نازوفارنکس؛ کنژنکتیویت در ۵۰\%، نومونی در ۲۰\%.
\item علائم کنژنکتیویت: \hl{ادم و ترشح ظرف ۵ تا ۱۴ روز تا چند هفته} پس از تولد؛ \hl{۱ تا ۲ هفته طول می‌کشد}.
\item درگیری تقریباً همیشه \hl{دوطرفه} است.
\item \hl{آدنوپاتی نمی‌دهد}؛ قرنیه را درگیر نمی‌کند.
\end{itemize}
\clearpage
\sheet{۹}
\subsection*{کلامیدیا ـ ادامه}
\begin{itemize}
\item نومونی: ظرف \hl{۲ تا ۱۹ هفته پس از تولد}؛ معمولاً \hl{بدون تب} است.
\begin{itemize}
\item رال، تاکی‌پنه، سرفه: \hl{\en{Repetitive staccato cough}}.
\item \hl{ویز ناشایع است}.
\item در \en{CXR} پرهوایی و درگیری منتشر دارد.
\item ممکن است در همراهی با \en{OM} و درگیری بینی باشد.
\end{itemize}
\item تشخیص: کشت نمونهٔ ملتحمه؛ \en{PCR}.
\item درمان: \hl{آزیترومایسین خوراکی} برای ۳ روز، یا \hl{اریترومایسین خوراکی} برای ۱۴ روز.
\begin{itemize}
\item \hl{درمان موضعی جایگاهی ندارد}.
\end{itemize}
\item پروفیلاکسی موضعی که برای همهٔ نوزادان انجام می‌شود، از کلامیدیا جلوگیری نمی‌کند.
\end{itemize}
\end{document}
